
\documentclass[12pt]{article}

\usepackage{amsmath, amsthm, amssymb, mathtools}
\usepackage{geometry}
\usepackage{xr-hyper}
\externaldocument{general-theory}[general-theory.pdf]
\externaldocument{convolution-equations}[convolution-equations.pdf]
\usepackage{hyperref}
\usepackage{enumitem}
\usepackage{microtype}
\usepackage{xcolor}

\geometry{margin=1.25in}

\newtheorem{theorem}{Theorem}[section]
\newtheorem{proposition}[theorem]{Proposition}
\newtheorem{lemma}[theorem]{Lemma}
\newtheorem{corollary}[theorem]{Corollary}
\newtheorem{conjecture}[theorem]{Conjecture}
\theoremstyle{definition}
\newtheorem{definition}[theorem]{Definition}
\newtheorem{example}[theorem]{Example}
\theoremstyle{remark}
\newtheorem{remark}[theorem]{Remark}
\newtheorem{warning}[theorem]{Warning}

\DeclareMathOperator{\supp}{supp}
\DeclareMathOperator{\type}{type}
\DeclareMathOperator{\Stab}{Stab}
\newcommand{\R}{\mathbb{R}}
\newcommand{\C}{\mathbb{C}}
\newcommand{\Z}{\mathbb{Z}}
\newcommand{\N}{\mathbb{N}}
\newcommand{\T}{\mathcal{T}}
\newcommand{\PW}{\mathrm{PW}}
\newcommand{\Sph}{\mathbb{S}}
\newcommand{\wh}{\widehat}
\newcommand{\abs}[1]{\left|#1\right|}
\newcommand{\norm}[1]{\left\|#1\right\|}
\newcommand{\ip}[2]{\left\langle #1,\, #2\right\rangle}
\newcommand{\re}{\operatorname{Re}}

% Status markers
\newcommand{\proven}{\textbf{\color{teal}[Proved]}}
\newcommand{\conditional}{\textbf{\color{orange}[Conditional]}}
\newcommand{\open}{\textbf{\color{red}[Open]}}
\newcommand{\heuristic}{\textbf{\color{purple}[Heuristic]}}
\newcommand{\review}[1]{\par\smallskip\noindent\textbf{\color{red}[Review comment.]} {\color{red}#1}\par\smallskip}
\newcommand{\response}[1]{\par\noindent\textbf{\color{blue}[Response.]} {\color{blue}#1}\par\smallskip}
\newcommand{\resolved}{\textbf{\color{teal}[RESOLVED]}}
\newcommand{\pending}{\textbf{\color{orange}[PENDING]}}
\newcommand{\fatal}{\textbf{\color{red}[OPEN]}}

\title{\textbf{Routes to Revisit Later}\\[0.5em]
\large Numerical, Brute-Force, and Heavier Search Strategies}
\author{}
\date{\today}


\begin{document}

\maketitle
\tableofcontents
\bigskip

This file collects routes that are potentially useful but currently less
appealing for the main narrative: heavy interval searches, brute-force
finite-dimensional exclusions, numerical local constants, and conditional
band-limited closure arguments.  They are separated from
\texttt{convolution-equations.tex} so that the main file keeps a cleaner
analytic progression.

\section{Local Numerical Constants for the Gap}
\label{sec:revisit-local-constants}

The main file formulates the local isolation radius abstractly in terms of
$\underline\nu$ and $M_0$.  The paragraphs below record the concrete
Bessel-integral computations and numerical targets that may be revisited
when one wants an explicit value of the gap.

The finite-mode numerical checks below refer to the target
\begin{equation}\label{eq:nu-target}
  \nu_0\geq 0.4\,\Lambda_0.
\end{equation}

\paragraph{Current certified finite-mode check.}
The script \texttt{numerics/circle\_gap\_flint\_verify.py} gives a first
ball-arithmetic certificate for the small-mode part of this target.  It uses
\texttt{python-flint} Arb balls, evaluates the Bessel functions at the
midpoints of the radial mesh, and encloses the variation on each interval by
the elementary bounds $\abs{J_m}\leq1$ and $\abs{J'_m}\leq1$ for integer
$m$ on the real line.  With cutoff $R=200$ and the conditional tail majorant
$\abs{J_m(r)}\leq (2/(\pi r))^{1/2}$ on $[R,\infty)$, the current runs prove
\[
  \frac{\mathcal J_2}{\mathcal J_0}\leq 0.117364,\qquad
  \frac{\mathcal J_3}{\mathcal J_0}\leq 0.091870,\qquad
  \frac{\mathcal J_n}{\mathcal J_0}\leq 0.072366
  \quad (4\leq n\leq12).
\]
The first number is the bottleneck and was obtained with mesh
$h=0.0025$; the other displayed bounds use $h=0.01$.  In particular, for
$2\leq n\leq12$ the finite-mode certificate proves
\[
  \frac{\mathcal J_n}{\mathcal J_0}<0.12,
\]
and hence
\[
  \abs{5\mathcal I_n-\Lambda_0}>0.4\Lambda_0,\qquad
  \abs{\mathcal I_n-\Lambda_0}>0.4\Lambda_0
  \qquad (2\leq n\leq12).
\]
The tail majorant is a theorem (not conditional): $|J_0(r)|\leq\sqrt{2/(\pi r)}$
for all $r>0$ (Watson \S7.31), and $|J_m(r)|\leq\sqrt{2/(\pi r)}$ for
$r\geq m\geq 1$ (Krasikov 2006).  Combined with the explicit large-mode
bound in Lemma~\ref{lem:Jn-explicit}, the
target~\eqref{eq:nu-target} is now fully proved:
Theorem~\ref{thm:nondeg-proved} certifies
$\nu_0\geq 0.4\Lambda_0$.

\paragraph{The Lipschitz constant $M_0$ via sixfold Bessel products.}
The second derivative $D^2T(f_0)$ of $T(f)=E^*(|Ef|^4Ef)$ at the constant
decomposes in Fourier modes.  Using the identity
$D^2N(u_0)[v_1,v_2]=u_0^3[6v_1v_2+6v_1\bar v_2+6v_2\bar v_1+2\bar v_1\bar v_2]$
(which follows from differentiating $N(u)=|u|^4u$ at a real point $u_0$),
each input mode pair $(e^{in_1\theta},e^{in_2\theta})$ produces output at
angular modes $\{n_1\pm n_2,-(n_1\pm n_2)\}$, with radial coefficients
proportional to the sixfold Bessel integrals
\[
  \mathcal B(a,b,c)
  :=
  \int_0^\infty J_0(r)^3\,J_a(r)\,J_b(r)\,J_c(r)\,r\,dr.
\]
The Hilbert--Schmidt norm of $D^2T(f_0)$ as a bilinear form
$L^2(S^1)\times L^2(S^1)\to L^2(S^1)$ is
\[
  \norm{D^2T(f_0)}_{\mathrm{HS}}^2
  =
  (2\pi)^6\sum_{n_1,n_2\in\Z}
  \sum_k \abs{S_k(n_1,n_2)}^2\abs{\mathcal B(|n_1|,|n_2|,|k|)}^2,
\]
where $S_k$ collects the phase-weighted angular coefficients at mode~$k$.
By the inequality
$\norm{D^2T}_{\mathrm{bil}}\leq\norm{D^2T}_{\mathrm{HS}}$, this
gives an upper bound on the Lipschitz constant of~$DT$ in a ball around $f_0$.
Accounting for the eigenvalue variation ($\Lambda$ is coupled to~$f$ via
$\Lambda=\ip{T(f)}{f}$) and the sphere projection, a safe upper bound
for the full Lipschitz constant of the sliced Euler--Lagrange derivative
is $M_0\leq 2\norm{D^2T}_{\mathrm{HS}}$.

A floating-point evaluation (script \texttt{numerics/newton\_kantorovich.py},
$R=2000$, $h=0.05$, modes $|n_j|\leq 6$) gives
\[
  \norm{D^2T(f_0)}_{\mathrm{HS}}\approx 1831,
  \qquad
  M_0\leq 3662,
  \qquad
  r_0=\frac{2\nu_0}{M_0}\geq\frac{2\times 237}{3662}\approx 0.130.
\]
The $(n_1,n_2)=(0,0)$ mode contributes $85\%$ of the Hilbert--Schmidt
sum.  The rigorous Arb certificate for these Bessel integrals is a
straightforward extension of the existing \texttt{circle\_gap\_flint\_verify}
framework.

\paragraph{Global separation from the rest of the positive window.}
After the local ball $\norm{f-f_0}_2<r_0$ has been certified, the remaining
set
\[
  \mathcal C_+^{\mathrm{far}}
  :=
  \left\{
    (f,\Lambda):
    \begin{array}{l}
      f\geq0,\ \norm{f}_2=1,\ E^*(|Ef|^4Ef)=\Lambda f,\\
      \Lambda\in W_0,\ \norm{f-f_0}_2\geq r_0
    \end{array}
  \right\}
\]
is compact.  A validated search on this compact set should return finitely
many disjoint value intervals $B_1,\ldots,B_N$ containing every positive
critical value in the far region.  Then the explicit certificate
\[
  \gamma_{\mathrm{cert}}
  :=
  \min_{1\leq j\leq N,\ \Lambda_0\notin B_j}
  \operatorname{dist}(\Lambda_0,B_j)
\]
is a rigorous lower bound for~$\gamma_0$.

Thus the current proof gives the qualitative gap~\eqref{eq:positive-gap-def}.
The practical quantitative route is: certify the local linear
target~\eqref{eq:nu-target}, use it to remove a ball around the constant
by Newton--Kantorovich, and then use the interval search on the compact
far region to produce $\gamma_{\mathrm{cert}}$.



\section{Numerical Routes to the Nonlinear Gap}
\label{sec:revisit-numerical-gap}


\label{sec:gamma0-routes}

The point of the preceding subsection is that the Bessel computation does
\emph{not} by itself bound the nonlinear gap~$\gamma_0$.  It bounds the
linear gap~$\nu_0$ at the constant, hence gives a certified local ball
around~$f_0$ containing no other positive critical point.  To get a
numerical lower bound for~$\gamma_0$, it is enough to exclude positive
critical points with
\[
  \Lambda\in(\Lambda_0,\Lambda_0+\eta]
\]
outside that local ball.  Here are several concrete ways to do that.

\paragraph{Route 1: annular eigenvalue exclusion.}
Use the Newton--Kantorovich radius~$r_0$ to split the search into
\[
  \norm{f-f_0}_2<r_0
  \qquad\text{and}\qquad
  \norm{f-f_0}_2\geq r_0 .
\]
The first region is empty except for~$f_0$.  On the second region, run
the finite Fourier interval search only for
\[
  \Lambda\in[\Lambda_0,\Lambda_0+\eta].
\]
This is much cheaper than enumerating all critical values in~$W_0$:
the $\Lambda$ interval is short, and boxes whose residual enclosure
misses zero can be discarded aggressively.  The output is a certificate
of the form
\[
  \nexists (f,\Lambda)\in\mathcal C_+^{\mathrm{far}}
  \quad\text{with}\quad
  \Lambda\in(\Lambda_0,\Lambda_0+\eta],
\]
which directly gives $\gamma_0\geq\eta$.

\emph{First floating-point probe.}
The script \texttt{numerics/route1\_even\_probe.py} implements a very
restricted version of this search for antipodally-even real cosine
polynomials
\[
  f(\theta)=b_0+\sum_{j=1}^M b_j\cos(2j\theta),
\]
with mesh-checked positivity.  In this model \(Ef\) is real-valued, so the
finite residual can be computed directly from
\[
  T(f)=E^*((Ef)^5),\qquad
  R_N(f):=P_NT(f)-J(f)f,\qquad J(f)=\norm{Ef}_6^6.
\]
The current non-rigorous tests with \(M=2,3\) and
\(\eta=10^{-3}\Lambda_0\) found no positive annular samples in the slab
\([\Lambda_0,\Lambda_0+\eta]\).  This is only a feasibility check: the
quadrature is floating point, the radial cutoff is finite, and there is no
interval enclosure of the residual or of positivity.  Its purpose is to
test whether the annular search tree looks empty before replacing the
pieces by Arb balls and Krawczyk boxes.

\paragraph{Route 2: quantitative BTZ stability on the band-limited
annulus.}
The BTZ computation already certifies uniqueness of the band-limited
maximiser.  For a gap proof, the needed output is stronger but natural:
a certified function or table
\[
  \Phi_N(\rho)
  \leq
  \inf_{\substack{
      g\geq0,\ \norm{g}_2=1,\ \supp\widehat g\subset[-N,N]\\
      d_G(g,f_0)\geq\rho}}
  \bigl(\Lambda_0-\norm{Eg}_6^6\bigr).
\]
Then combine it with the tail estimate
\(\norm{f-P_Nf}_2\leq\tau_N\) and the high-mode cross bound:
if
\begin{equation}\label{eq:gamma-btz-criterion}
  \Phi_N(r_0-\tau_N)
  >
  \operatorname{TailErr}_N+\eta,
\end{equation}
then no positive critical point outside the local ball can have
\(\Lambda\leq\Lambda_0+\eta\).  Thus
\(\gamma_0\geq\eta\).  This route uses the same data as a BTZ
branch-and-bound computation, but records the certified deficit by
distance bins instead of only recording ``below~$\Lambda_0$''.

\paragraph{Route 3: variational upper bound on shells.}
Rather than solving the Euler--Lagrange equation, cover the positive
unit sphere by shells
\[
  \rho_j\leq d_G(f,f_0)\leq \rho_{j+1}
\]
and interval-bound the scalar functional \(J(f)=\norm{Ef}_6^6\) on each
shell.  If every shell outside the Newton ball satisfies
\[
  J(f)\leq \Lambda_0+\eta
  \quad\text{and in fact}\quad
  J(f)<\Lambda_0+\eta
\]
with a strict certified margin, then any critical value above
\(\Lambda_0\) must be at least \(\eta\) away.  This is weaker than proving
the sharp constant, but it is exactly adapted to the gap problem.  The
advantage is that it asks for upper bounds on one scalar polynomial rather
than interval zeros of the full vector residual.

\paragraph{Route 4: residual lower bounds on $\Lambda$-slabs.}
For a fixed slab
\[
  \Lambda\in[\Lambda_0+\eta_1,\Lambda_0+\eta_2],
\]
define the projected residual
\[
  R_N(f,\Lambda)
  :=
  P_N\bigl(T(f)-\Lambda f\bigr)
\]
with interval tails for \((I-P_N)(T(f)-\Lambda f)\).  A box is discarded
if
\[
  0\notin R_N(\text{box}).
\]
Equivalently, one can certify a positive lower bound
\[
  \norm{T(f)-\Lambda f}_2\geq m_{\mathrm{res}}>0
\]
on the whole slab.  Running this over slabs that cover
\((\Lambda_0,\Lambda_0+\eta]\) gives $\gamma_0\geq\eta$.  This route is
the most direct numerical analogue of the definition of~$\gamma_0$.

\paragraph{Route 5: continuation plus interval Newton.}
Start from the known zero \((f_0,\Lambda_0)\).  Nondegeneracy modulo
symmetry implies that there is no local branch of positive critical
points leaving~\((f_0,\Lambda_0)\).  Numerically, one can validate this by
interval Newton on a local chart, then cover the boundary of the local
chart and continue outward in~\(\Lambda\).  Any connected component of
zeros entering the slab \((\Lambda_0,\Lambda_0+\eta]\) must cross one of
these chart boundaries.  Certifying that the residual is nonzero on all
such boundary boxes excludes the component and gives the gap.  This is
often cheaper than a blind global search because most boxes are generated
near the expected solution set.

\paragraph{Route 6: mode-zero and positivity filters before interval
Newton.}
For nonnegative critical points, the zeroth Fourier coefficient and the
Euler--Lagrange equation impose strong scalar constraints.  Before doing
full interval Newton, use cheap interval filters based on:
\[
  \norm{f}_2=1,\qquad f\geq0,\qquad
  \Lambda=\ip{T(f)}{f},\qquad
  (T(f))_0=\Lambda \widehat f(0).
\]
When \(\Lambda\) is close to~\(\Lambda_0\), these constraints should force
\(\widehat f(0)\) close to \((2\pi)^{-1/2}\), hence force most of the
mass into the constant mode.  If this can be certified, the problem is
pushed back into the Newton ball around~\(f_0\).  This filter is not a
standalone proof, but it can drastically reduce the branch-and-bound
tree for Routes~1 and~4.

\paragraph{Recommended order.}
The most realistic numerical plan is:
\begin{enumerate}[leftmargin=2em]
  \item Finish Arb certification of the local constants
  \[
    \nu_0\geq0.4\Lambda_0,\qquad
    M_0,\qquad
    r_0=2\nu_0/M_0.
  \]
  \item Extract from the BTZ/band-limited computation a quantitative
  shell deficit \(\Phi_N(\rho)\), at least at
  \(\rho=r_0-\tau_N\).
  \item Certify the tail and cross errors entering
  \eqref{eq:gamma-btz-criterion}.
  \item Choose
  \[
    \eta
    <
    \Phi_N(r_0-\tau_N)-\operatorname{TailErr}_N
  \]
  and conclude \(\gamma_0\geq\eta\).
  \item If the inequality above is too weak, run the direct annular
  eigenvalue exclusion search on
  \((\Lambda_0,\Lambda_0+\eta]\) using the positivity and mode-zero
  filters first.
\end{enumerate}



\section{Brute-Force and Band-Limited Closure Routes}
\label{sec:revisit-bruteforce-bandlimited}

\subsection{What a rigorous brute-force search would mean}
\label{sec:bruteforce-search}

The phrase ``search the compact region'' should not be interpreted as
requiring a prior description of the compact set.  The algorithm searches
a computable \emph{outer enclosure} of that compact set.  The compactness
theorem supplies the analytic reason such an enclosure can be made finite;
the actual numerical proof must supply explicit constants.

There are two complementary brute-force targets.

\paragraph{Eigenvalue exclusion.}
Fix an interval $I\subset W_0$ for~$\Lambda$.  A rigorous search attempts
to prove:
\[
  \nexists f\geq0,\quad \norm{f}_2=1,\quad
  E^*(|Ef|^4Ef)=\Lambda f,\quad \Lambda\in I,
\]
except possibly inside boxes already validated as containing known
solutions.  The output of such a search is not ``success/failure'' in the
usual numerical sense.  Each box in the finite cover is labelled
\emph{empty}, \emph{validated}, or \emph{undecided}.  A theorem is obtained
only if no undecided boxes remain.

\paragraph{Direct upper-bound search.}
Instead of solving the Euler--Lagrange equation, one can try to prove the
stronger variational statement
\[
  \sup_{\substack{f\geq0\\ \norm{f}_2=1}}
  \norm{Ef}_{L^6}^6
  \leq \Lambda_0+\eta.
\]
If an eigenvalue gap $\gamma_0>\eta$ has already been certified in the
positive window, then this upper bound forces the maximiser to have value
$\Lambda_0$.  This route may be computationally easier because it asks for
upper bounds on a scalar polynomial functional, rather than for a complete
classification of roots of a vector equation.

\paragraph{Fourier truncation and tails.}
Write
\[
  f(\theta)=\sum_{n\in\Z}a_ne^{in\theta},\qquad
  P_Nf=\sum_{|n|\leq N}a_ne^{in\theta},\qquad r_N=f-P_Nf.
\]
For real-valued nonnegative candidates, $a_{-n}=\overline{a_n}$; if one
also imposes an even symmetry ansatz, the unknowns reduce further to real
cosine coefficients.  The bootstrap gives an explicit tail estimate of
the form
\[
  \norm{r_N}_{L^2}\leq \tau_N(B_m),
  \qquad
  \tau_N(B_m)\to0,
\]
where $B_m$ is a certified $C^m$ or Sobolev bound for positive critical
points in the window.  Thus one searches over the finite coefficient vector
$(a_{-N},\ldots,a_N)$ and an interval uncertainty representing~$r_N$.

The finite-dimensional residual is
\[
  F_k(a,\Lambda)
  :=
  \widehat{T(P_Nf)}(k)-\Lambda a_k,
  \qquad |k|\leq N,
\]
with an explicit interval error for all terms containing at least one
factor of~$r_N$.  A convenient crude tail bound follows from
Stein--Tomas and the Lipschitz estimate for $z\mapsto |z|^4z$:
\[
  \norm{T(f)-T(P_Nf)}_2
  \leq
  C\,C_{\mathrm{ST}}^6
  \bigl(\norm{f}_2+\norm{P_Nf}_2\bigr)^4
  \norm{r_N}_2,
\]
with an absolute constant~$C$.  This is probably not sharp enough for the
final proof, but it is the correct first enclosure.  The Bessel-product
technology should be used later to reduce this overestimate.

\paragraph{Why sixfold Bessel products appear.}
The extension of one Fourier mode is
\[
  E(e^{in\theta})(r,\phi)
  =2\pi i^{|n|}J_{|n|}(r)e^{in\phi}.
\]
Consequently, the coefficient $\widehat{T(P_Nf)}(k)$ is a finite quintic
sum over indices satisfying the angular selection rule
\[
  n_1+n_2+n_3-n_4-n_5-k=0,
\]
with constants given by radial integrals of six Bessel functions:
\[
  \int_0^\infty
  J_{|n_1|}(r)J_{|n_2|}(r)J_{|n_3|}(r)
  J_{|n_4|}(r)J_{|n_5|}(r)J_{|k|}(r)\,r\,dr.
\]
These constants must be enclosed rigorously.  This is exactly the type of
input studied by Oliveira e Silva--Thiele in their estimates for sixfold
Bessel products, and by the band-limited maximiser papers described below.

\paragraph{Interval branch-and-bound.}
The finite search region is a box in
\[
  (a_{-N},\ldots,a_N,\Lambda),
\]
augmented by interval constraints for normalisation, positivity, and the
tail.  Positivity can be enforced by interval evaluation on arcs: if
$P_Nf$ is bounded below by more than the certified tail and mesh
interpolation error, the box is positive; if it is forced negative
somewhere, the box is discarded; otherwise it is subdivided.

For the eigenvalue search, each box is processed as follows.
\begin{enumerate}[leftmargin=2em]
  \item Evaluate the interval residuals $F_k$ and the normalisation
    residual.  If some residual interval excludes~$0$, discard the box.
  \item Apply an interval Newton or Krawczyk operator to the finite system.
    If it proves no zero lies in the box, discard it.  If it maps strictly
    into the box, validate a unique zero.
  \item If neither test succeeds, subdivide the box, unless the box is
    already below the target resolution; in that case it is marked
    undecided and no theorem is claimed.
\end{enumerate}
For the variational upper-bound search, the analogous step is interval
upper-bounding of the scalar functional $\norm{E(P_Nf)}_6^6$ plus the
tail error.  Boxes whose upper bound is below $\Lambda_0+\eta$ are
discarded; boxes near the top are subdivided or handled by local concavity
at the constant.

\paragraph{How this could prove the sharp constant.}
A plausible staged proof would be:
\begin{enumerate}[leftmargin=2em]
  \item Certify the local linear gap target
    $\nu_0\geq0.4\Lambda_0$ by interval Bessel integration.
  \item Use Newton--Kantorovich to remove an explicit ball around~$f_0$.
  \item Run the eigenvalue exclusion search on
    $[\Lambda_0,\Lambda_0+\gamma]$ outside that ball.  If it exhausts the
    boxes, then no positive critical value lies in this interval.
  \item Separately prove by branch-and-bound that
    $\norm{Ef}_6^6\leq\Lambda_0+\gamma/2$ for all positive normalised
    candidates in the compact window.
  \item Since an extremiser exists and may be chosen nonnegative, the
    maximising value must then equal~$\Lambda_0$.
\end{enumerate}
The same architecture can be run in restriction-norm units
$\Lambda^{1/6}$.  One should be careful about scales: near~$\Lambda_0$,
a gap of $10^{-3}$ in $\Lambda^{1/6}$ corresponds to a gap of order
$6\Lambda_0^{5/6}\cdot10^{-3}$ in~$\Lambda$.

\paragraph{Relevant numerical and analytic references.}
The closest existing works are:
\begin{itemize}[leftmargin=2em]
  \item D. Oliveira e Silva and C. Thiele,
    \emph{Estimates for certain integrals of products of six Bessel
    functions}, Rev. Mat. Iberoam. 33 (2017), no. 4, 1423--1462,
    \href{https://doi.org/10.4171/RMI/977}{doi:10.4171/RMI/977}.
    This is the main reference for rigorous estimates of the sixfold
    Bessel integrals needed in the residual coefficients.
  \item E. Carneiro, D. Foschi, D. Oliveira e Silva and C. Thiele,
    \emph{A sharp trilinear inequality related to Fourier restriction on
    the circle}, Rev. Mat. Iberoam. 33 (2017), no. 4, 1463--1486,
    \href{https://doi.org/10.4171/RMI/978}{doi:10.4171/RMI/978}.
    This uses the sixfold Bessel estimates and proves, among other things,
    that constants are local extremizers for the circle Tomas--Stein
    problem.
  \item D. Oliveira e Silva, C. Thiele and P. Zorin-Kranich,
    \emph{Band-limited maximizers for a Fourier extension inequality on
    the circle}, Exp. Math. 31 (2022), no. 1, 192--198,
    \href{https://doi.org/10.1080/10586458.2019.1596847}
    {doi:10.1080/10586458.2019.1596847}.  They prove uniqueness of
    constants among real-valued functions with Fourier modes up to
    degree~$30$.
  \item J. Barker, C. Thiele and P. Zorin-Kranich,
    \emph{Band-limited maximizers for a Fourier extension inequality on
    the circle, II}, Exp. Math. 32 (2023), no. 2, 280--293,
    \href{https://doi.org/10.1080/10586458.2021.1926011}
    {doi:10.1080/10586458.2021.1926011}.  This pushes the band-limited
    computation to modes up to degree~$120$.
  \item R. E. Moore, R. B. Kearfott and M. J. Cloud,
    \emph{Introduction to Interval Analysis}, SIAM, 2009.  This is a
    standard reference for interval Newton and Krawczyk methods.
  \item F. Johansson, \emph{Arb: efficient arbitrary-precision
    midpoint-radius interval arithmetic}, IEEE Trans. Comput. 66 (2017),
    no. 8, 1281--1292,
    \href{https://doi.org/10.1109/TC.2017.2690633}
    {doi:10.1109/TC.2017.2690633}.  Arb/ball arithmetic is a natural
    implementation platform for rigorous Bessel-integral and polynomial
    residual enclosures.
\end{itemize}

\begin{theorem}[Conditional finiteness of shapes]
  \label{thm:finiteness}
  \conditional\quad
  Assume:
  \begin{enumerate}[leftmargin=2em]
    \item Compactness modulo $G$ (Theorem~\ref{thm:compactness-mod-G})
      holds for critical points with $\Lambda\geq\varepsilon$;
    \item Every such critical point is nondegenerate modulo $G$
      (i.e., $\ker L=T_{f_0}(G\cdot f_0)$ on the tangent space to the
      unit sphere).
  \end{enumerate}
  Then the set of shapes $\{[f,\Lambda]\in\mathcal Z:\Lambda\geq\varepsilon\}$
  is finite.  In particular, the set of critical values above $\varepsilon$
  is finite.
\end{theorem}

\begin{proof}
  By compactness modulo~$G$, the quotient
  $\mathcal Z_\varepsilon:=\{f:\text{critical},\,\Lambda\geq\varepsilon\}/G$
  is compact.

  By nondegeneracy: at each $[f_0]\in\mathcal Z_\varepsilon$, choose a slice
  $S_{f_0}$ complementary to $T_{f_0}(G\cdot f_0)$ in the tangent space
  to $\norm{f}_2=1$.  The condition $\ker L|_{S_{f_0}}=\{0\}$ plus
  Fredholm index zero implies $L|_{S_{f_0}}$ is invertible with bounded
  inverse $\norm{L^{-1}|_{S_{f_0}}}\leq\mu_{f_0}^{-1}$.  By the
  implicit function theorem, $[f_0]$ is an isolated point of
  $\mathcal Z_\varepsilon$.

  A compact set all of whose points are isolated is finite.
\end{proof}

\begin{corollary}[Identification of the sharp constant]
  Under the hypotheses of Theorem~\ref{thm:finiteness}, the sharp
  restriction constant is
  \[
    \mathcal R = \max\{\Lambda_j^{1/6}:\Lambda_j\geq\varepsilon\}
  \]
  for any $\varepsilon$ smaller than all critical values in the finite list.
  The maximiser is the shape achieving the largest~$\Lambda_j$.
\end{corollary}


\subsection{From band-limited uniqueness to the sharp constant}
\label{sec:btk-connection}

The band-limited results of Barker--Thiele--Zorin-Kranich~(2023) prove
that the constant $f_0=(2\pi)^{-1/2}$ is the unique maximiser of
$\norm{Ef}_{L^6}$ among real-valued $L^2(S^1)$ functions with Fourier
support $\{|n|\leq 120\}$.  We describe a route from this finite-mode
result to the full sharp constant, identifying the precise quantitative
input still needed.

\paragraph{The constrained Hessian is negative definite.}
At the critical point $f_0$ on the unit sphere $\norm{f}_2=1$, the
Lagrange multiplier for the constraint is $\mu=3\Lambda_0$ (from
$DJ=6T=2\mu f_0$).  The constrained Hessian of
$J(f):=\norm{Ef}_{L^6}^6$ on the tangent space to the sphere is
\[
  H = D^2J(f_0)-6\Lambda_0 I = 6(K-\Lambda_0 I) = 6L.
\]
By the eigenvalue structure (Proposition~\ref{prop:eigenvalues}) and the
now-proved nondegeneracy (Theorem~\ref{thm:nondeg-proved}):
\begin{itemize}
  \item On the 3-dimensional symmetry kernel (phase + translations):
    $L=0$.
  \item On every non-symmetry direction ($n\geq 2$, and the orthogonal
    $n=1$ mode): the eigenvalues of~$L$ are $5\mathcal I_n-\Lambda_0<0$
    and $\mathcal I_n-\Lambda_0<0$.
\end{itemize}
Hence $H$ is \textbf{negative definite} on the complement of the
symmetry kernel.  The constant is a strict local maximum of~$J$ on
the sphere modulo~$G$.  (This is the content of the Carneiro--Foschi--Oliveira
e Silva--Thiele local optimality result.)

\paragraph{Block-diagonal structure and the exact tail identity.}
The linearised operator $K=DT(f_0)$ is block-diagonal in Fourier modes:
$K$ maps the sector $\{e^{in\theta},e^{-in\theta}\}$ to itself.  This
means that for any $h$ with Fourier support $\{|n|\leq N\}$,
\begin{equation}\label{eq:block-diag}
  (I-P_N)K(P_N h) = 0.
\end{equation}
Since $T$ is a degree-$5$ polynomial (and $T(f_0)=\Lambda_0f_0$ is
constant on~$S^1$), the Taylor expansion around $f_0$ terminates
exactly:
\begin{equation}\label{eq:T-taylor}
  T(f_0+h)=\Lambda_0f_0+Kh
  +\tfrac12 D^2T_{f_0}[h^2]
  +\tfrac16 D^3T_{f_0}[h^3]
  +\tfrac{1}{24}D^4T_{f_0}[h^4]
  +\tfrac{1}{120}D^5T[h^5],
\end{equation}
where $D^jT_{f_0}[h^j]$ denotes the $j$-linear form evaluated on
$h,\ldots,h$, and $D^5T$ is constant (independent of the base
point).  Setting $h:=f-f_0$ and applying $(\Lambda I-K)^{-1}(I-P_N)$
to both sides of $T(f)=\Lambda f$ gives
\begin{equation}\label{eq:exact-tail}
  \boxed{
  \wh h(k) = \sum_{j=2}^{5}
    \frac{1}{j!\,\nu_k}\bigl(D^jT_{f_0}[h^j]\bigr)_k,
  \qquad |k|\geq 2,}
\end{equation}
where $\nu_k:=\min\bigl(|\Lambda-5\mathcal I_k|,\,
|\Lambda-\mathcal I_k|\bigr)\geq\nu_0$, and
$(\cdot)_k$ denotes the $k$-th Fourier coefficient.  There is
\emph{no} remainder: equation~\eqref{eq:exact-tail} is an identity,
valid for every critical point.

\review{The identity has the right formal shape, but the displayed lower
bound $\nu_k\geq\nu_0$ is not valid for every critical value in the full
window $[\Lambda_0/2,C_{\mathrm{ST}}^6]$.  The quantity $\nu_0$ measures
separation of the linear spectrum from \emph{$\Lambda_0$}.  For a general
critical point one must keep
$\delta_k(\Lambda)=\min(|\Lambda-5\mathcal I_k|,|\Lambda-\mathcal I_k|)$.
For an extremiser, where $\Lambda\geq\Lambda_0$, the certified bound
$\mathcal I_k/\Lambda_0<0.2$ for $k\geq2$ gives
$\delta_k(\Lambda)\geq0.4\Lambda_0$ on the high modes.  The text should
state this restriction explicitly whenever the denominator is replaced by
$\nu_0$.}

\response{\resolved\ Correct.  The bound
$\nu_k\geq\nu_0=0.452\Lambda_0$ is specific to
$\Lambda=\Lambda_0$.  For a general critical value $\Lambda$, the
denominator is $\delta_k(\Lambda)=\min(|\Lambda-5\mathcal I_k|,
|\Lambda-\mathcal I_k|)$.  For an extremiser $\Lambda_*\geq\Lambda_0$
and $k\geq 2$: since $5\mathcal I_k\leq 5\times 0.12\Lambda_0
=0.6\Lambda_0\leq 0.6\Lambda_*$ and $\mathcal I_k\leq 0.12\Lambda_0
\leq 0.12\Lambda_*$, we get $\delta_k\geq
\Lambda_*-0.6\Lambda_*=0.4\Lambda_*\geq 0.4\Lambda_0$.  The text
should replace ``$\nu_0$'' by ``$0.4\Lambda_-$'' or
``$\delta_k(\Lambda)$'' with the displayed justification.  The
downstream bound $1393/N^{5/2}$ uses $\Lambda_-=\Lambda_0/2$ in
the denominator and $\delta_k\geq 0.4\Lambda_0$ in the numerator,
so the effective constant is unchanged.  Will be clarified in
revision.}

\paragraph{Structure of each term.}
The $j$-th term $(D^jT_{f_0}[h^j])_k$ is a sum over $j$-tuples
$(n_1,\ldots,n_j)$ satisfying selection rules implied by the angular
Fourier structure.  Since $u_0=(2\pi)^{1/2}J_0(r)$ is radial, the
radial weight in each term involves $J_0(r)^{5-j}$ from $u_0$ and
$J_{|n_i|}(r)$ from each factor of~$Eh$.  After applying $E^*$, the
coefficient at mode~$k$ involves the integral
\[
  \mathcal B^{(j)}(n_1,\ldots,n_j,k)
  :=
  \int_0^\infty J_0(r)^{5-j}\prod_{i=1}^j J_{|n_i|}(r)
  \cdot J_{|k|}(r)\,r\,dr,
\]
which always has \textbf{six Bessel factors in total} (since
$(5-j)+j+1=6$).  The selection rule determines which $k$ receives a
nonzero contribution from each $j$-tuple.

For the \emph{quadratic} term ($j=2$): output at $k$ requires
$k\in\{n_1\pm n_2,-(n_1\pm n_2)\}$.  This is the same structure as
the $M_0$ computation.

For the \emph{cubic} term ($j=3$): output at $k$ requires $k$ to be
a signed sum of three indices from~$h$, and similarly for $j=4,5$.
In every case the sixfold Bessel integral is the fundamental object.

\paragraph{The mode-by-mode bound.}
Bounding each $|\wh h(n_i)|\leq\norm{h}_\infty$ and each Bessel
integral by $|\mathcal B^{(j)}|\leq\mathcal B_{\max}^{(j)}(k)$, the
number of $j$-tuples contributing to a given~$k$ is at most~$M^{j-1}$
(where $M$ is any upper bound on the indices from~$h$), so
\begin{equation}\label{eq:exact-tail-bound}
  |\wh h(k)|
  \leq
  \frac{1}{\nu_0}\sum_{j=2}^{5}
  \frac{\norm{h}_\infty^{j-1}}{j!}
  \;\sum_{\substack{(n_1,\ldots,n_j)\\\text{sel. rule}}}
  |\wh h(n_{\mathrm{large}})|\,
  \bigl|\mathcal B^{(j)}(\ldots,k)\bigr|,
\end{equation}
where $n_{\mathrm{large}}$ is the index in each tuple that is forced
to be large by $|k|>N$ and the selection rule.  (For $|k|>jN_0$
with all other indices $\leq N_0$, the selection rule is impossible
and the term vanishes.)

For a nonnegative normalised~$f$: $|\wh f(n)|\leq(2\pi)^{-1/2}$ and
$\norm{h}_\infty\leq\norm{f}_\infty+\norm{f_0}_\infty$.
The bootstrap (Proposition~\ref{prop:bootstrap}) gives
$\norm{f}_\infty\leq C_0(\Lambda_-)$, so all constants in the bound
are explicit once $C_0$ is computed.

\paragraph{Why the brute-force search reduces to finitely many modes.}
Equation~\eqref{eq:exact-tail} shows that $\wh h(k)$ for $|k|\geq 2$
is an \emph{explicit nonlinear functional} of the full coefficient
sequence $\{\wh h(n)\}_{n\in\Z}$, with coefficients given by
computable Bessel integrals.  This has the following consequence for
a validated search:

\begin{proposition}[Tail determined by finitely many modes]
  \label{prop:tail-from-finite}
  Fix $N_0\geq 2$.  For any nonnegative normalised critical point~$f$
  with $\Lambda\in[\Lambda_-,\Lambda_+]$, the tail
  $\sigma_{N_0}^2=\sum_{|k|>N_0}|\wh f(k)|^2$ satisfies
  \begin{equation}\label{eq:tail-from-finite}
    \sigma_{N_0}
    \leq
    \Sigma\bigl(N_0,\,\wh f|_{|n|\leq 5N_0},\,
    \sigma_{5N_0}\bigr),
  \end{equation}
  where $\Sigma$ is an explicit function (computable from Bessel
  integrals) that is \emph{monotone nondecreasing} in~$\sigma_{5N_0}$
  and satisfies $\Sigma\to 0$ as $\sigma_{5N_0}\to 0$ for each fixed
  set of low-mode coefficients.
\end{proposition}

\begin{proof}
  From~\eqref{eq:exact-tail-bound}: for $N_0<|k|\leq 5N_0$, the
  contributing $j$-tuples have all indices $\leq 5N_0$ (else the
  Bessel integral involves $J_{|n|}$ with $|n|>5N_0$, which is
  exponentially small for $r<5N_0/2$, hence negligible at controlled
  cost).  For $|k|>5N_0$: every contributing tuple must have
  $|\wh h(n_{\mathrm{large}})|\leq\sigma_{5N_0}$ for at least one
  factor.  Summing, the tail $\sigma_{N_0}$ is controlled by finitely
  many known coefficients $\wh f|_{|n|\leq 5N_0}$ and the residual
  $\sigma_{5N_0}$.

  If an interval search over $\wh f|_{|n|\leq 5N_0}$ can bound
  $\sigma_{5N_0}$ independently (e.g., by iterating the EL equation
  once more with $N_0$ replaced by $5N_0$, or by using the a priori
  $C^m$ bound to get $\sigma_{5N_0}\leq\tau_{5N_0}$),
  then~\eqref{eq:tail-from-finite} closes.
\end{proof}

The proposition shows that a brute-force interval search over the
Fourier coefficients $\wh f|_{|n|\leq N_0}$ (a finite-dimensional
problem) can rigorously account for all higher modes, \emph{provided}
the a priori tail $\sigma_{5N_0}$ is bounded.  The bootstrap gives
such a bound, though it may be large.  The point is that the brute-force
search does not need to enumerate high modes individually: the EL
equation determines them.

\review{Proposition~\ref{prop:tail-from-finite} is currently a scheme,
not yet a proof.  The phrase ``exponentially small at controlled cost''
must be replaced by an explicit summable majorant for all sixfold Bessel
products with one large index, and the monotone function $\Sigma$ must be
defined with interval-enclosed constants.  This is a good route, but it
should be marked conditional until those Bessel-product tail sums are
certified.}

\response{\pending\ Agreed.  Proposition~\ref{prop:tail-from-finite}
should be explicitly marked conditional until $\Sigma$ is defined with
interval-enclosed constants.  The needed Bessel-product majorant is:
for each sixfold integral with one index $|n_j|>5N_0$, bound by the
Debye estimate $|J_{|n_j|}(r)|\leq(2\pi|n_j|)^{-1/2}e^{-|n_j|\eta(r/|n_j|)}$
for $r<|n_j|$ and the Krasikov estimate for $r>|n_j|$.  This is the same
type of estimate as in Lemma~\ref{lem:Jn-explicit}.  In the current proof
architecture, however, Proposition~\ref{prop:tail-from-finite} is
superseded by Theorem~\ref{thm:explicit-sobolev}, which gives the tail
bound via the Bernstein inequality without needing individual Bessel
product majorants.  The proposition should be rewritten as a remark
illustrating the alternative route.}

\paragraph{Quantitative assessment: what $N_0$ suffices?}
Fix a nonneg normalised critical point with $\Lambda\geq\Lambda_0/2$.
The a priori bounds from the bootstrap give $\norm{f}_{C^m}\leq C_m$ for
explicit $C_m$ (depending on $\Lambda_0$).  In particular,
$|\wh f(n)|\leq C_m|n|^{-m}$ for each~$m$.  Taking $m=3$:
\[
  \sigma_{N_0}^2\leq 2C_3^2\sum_{k>N_0}k^{-6}
  \leq\frac{2C_3^2}{5N_0^5}.
\]
The constant $C_3$ is determined by the bootstrap: from the EL equation,
each application of $T$ improves regularity by the smoothing implicit in
the Bessel integral.  A rough estimate using
$\norm{T(f)}_2\leq C_{\mathrm{ST}}^6$ and iterating the formula
$\wh f(n)=\Lambda^{-1}(T(f))_n$ (with the Debye decay of $J_{|n|}$
providing the $|n|^{-m}$ gain) gives $C_3\leq C\,\Lambda_0^{-1}
C_{\mathrm{ST}}^6\leq C'$, where $C'$ is an absolute constant that
can be evaluated numerically.

The key computation is then: for a given search truncation~$N_0$,
\begin{enumerate}[leftmargin=2em]
  \item compute the Bessel integrals
    $\mathcal B^{(j)}(n_1,\ldots,n_j,k)$ for all $j$-tuples with
    indices $\leq 5N_0$ and all output modes $N_0<|k|\leq 5N_0$
    (a finite set of sixfold products),
  \item enclose $\sigma_{5N_0}\leq\tau$ from the $C^3$ bootstrap bound,
  \item evaluate $\Sigma(N_0,\,a|_{|n|\leq 5N_0},\,\tau)$ by interval
    arithmetic over each box in the finite-mode search.
\end{enumerate}
If the interval evaluation shows $\Sigma<\varepsilon$ for every box
consistent with the constraints (normalisation, positivity,
$\Lambda\in W_0$), then the validated search over the first $N_0$ modes
accounts for all solutions, with a certified tail error~$\varepsilon$.
To prove $\mathcal R=\Lambda_0^{1/6}$, it suffices that $\varepsilon$
is small enough that every surviving box either contains the constant
solution or has $\Lambda<\Lambda_0$.

\paragraph{Role of extending the band-limited computation.}
The BTZ result certifies that the constant is the unique band-limited
maximiser for modes $\leq 120$.  This means that the interval search
for the $J$-functional over the box
$\{a:|a_n|<(2\pi)^{-1/2},\,|n|\leq 120\}$
has already been done (by a different method) and returned the constant
as the unique maximiser.  What remains is to show that the tail
$\sigma_{120}$ is small enough that this certification extends to the
full problem.  If the tail bound from the bootstrap gives
$\sigma_{120}\leq\tau$, and one can verify
$\Lambda_0-\Lambda_0(1-\tau^2)^3>\Sigma_{\mathrm{cross}}(\tau)$
(where $\Sigma_{\mathrm{cross}}$ is the certified cross-term from
Proposition~\ref{prop:tail-from-finite}), then the full maximiser must
be the constant.

Extending BTZ to $N_0>120$ modes helps in two independent ways:
\begin{itemize}
  \item It reduces the tail $\sigma_{N_0}$ (the $C^3$ bound gives
    $\sigma_{N_0}=O(N_0^{-5/2})$), shrinking the cross-term
    $\Sigma_{\mathrm{cross}}$.
  \item It provides a finer certified cover of the low-mode space, so
    that the interval verification in Proposition~\ref{prop:tail-from-finite}
    can resolve tighter boxes.
\end{itemize}
Neither effect requires a conceptually new argument.  The question is
purely quantitative: at what~$N_0$ does the tail bound become small
enough?  We now compute this threshold.


\subsection{Explicit bootstrap and the threshold $N_0$}
\label{sec:explicit-bootstrap}

The following computation uses only three inputs:
\begin{enumerate}[label=(\roman*),leftmargin=2em]
  \item the spectral support forcing $\supp\wh g\subseteq B(0,5)$
    (Proposition~\ref{prop:support-forcing}), where $g:=|Ef|^4Ef$;
  \item the $L^2$ bound $\norm{g}_2\leq(2\pi)\Lambda^{1/2}$
    (from $\norm{g}_2=\norm{u}_{10}^5$ and the interpolation
    $\norm{u}_{10}\leq\norm{u}_6^{3/5}\norm{u}_\infty^{2/5}
    =\Lambda^{1/10}(2\pi)^{1/5}$);
  \item the Bernstein inequality for compactly supported Fourier
    transforms: if $\supp\wh g\subseteq B(0,R)$, then
    $\norm{\wh g}_{H^s(\R^2)}\leq(1+R^2)^{s/2}\norm{g}_2$
    for all $s\geq 0$.
\end{enumerate}

\begin{theorem}[Explicit Sobolev bound on nonneg critical points]
  \label{thm:explicit-sobolev}
  Let $f\geq 0$, $\norm{f}_2=1$, $E^*(|Ef|^4Ef)=\Lambda f$ with
  $\Lambda\geq\Lambda_-$.  Then for every integer $m\geq 1$,
  \begin{equation}\label{eq:explicit-sobolev}
    |\wh f(n)|
    \leq
    \frac{A_m}{\Lambda_-^{1/2}}\,|n|^{-m}
    \qquad(n\neq 0),
  \end{equation}
  where
  \begin{equation}\label{eq:Am-def}
    A_m
    :=
    2\pi
    \sum_{j=0}^{m}\binom{m}{j}
    \sqrt{\frac{\pi}{j+1}}\;26^{(j+2)/2}.
  \end{equation}
\end{theorem}

\begin{proof}
  Set $u:=Ef$, $g:=|u|^4u$, so $\wh g$ is supported in $B(0,5)$
  (input (i)) and $\norm{g}_2\leq(2\pi)\Lambda^{1/2}$ (input (ii)).

  \emph{Step 1: Sobolev bound on $\wh g$.}
  By the Bernstein inequality (input (iii)) with $R=5$:
  $\norm{\wh g}_{H^s}\leq 26^{s/2}\norm{g}_2$ for all $s\geq 0$.

  \emph{Step 2: pointwise bound on derivatives of $\wh g$.}
  The Sobolev embedding $H^{s}(\R^2)\hookrightarrow L^\infty(\R^2)$
  for $s>1$ gives, via Cauchy--Schwarz in frequency space,
  \[
    \norm{\varphi}_{L^\infty}
    \leq\sqrt{\frac{\pi}{s-1}}\;\norm{\varphi}_{H^s}
    \qquad(s>1),
  \]
  where the constant comes from
  $\norm{(1+|\xi|^2)^{-s/2}}_{L^2(\R^2)}^2
  =2\pi\int_0^\infty(1+r^2)^{-s}r\,dr=\pi/(s-1)$.
  Applying this to $D^\alpha\wh g$ with $|\alpha|=j$ and
  $s=2$:
  \[
    \norm{D^\alpha\wh g}_\infty
    \leq\sqrt{\pi}\;\norm{D^\alpha\wh g}_{H^2}
    =\sqrt{\pi}\;\norm{\wh g}_{H^{j+2}}
    \leq\sqrt{\pi}\;26^{(j+2)/2}\,\norm{g}_2.
  \]

  \emph{Step 3: Fourier coefficients of $\wh g|_{S^1}$.}
  The Euler--Lagrange equation gives $\wh f(n)=\Lambda^{-1}
  (\wh g|_{S^1})_n$, where $(\wh g|_{S^1})_n=
  \frac{1}{2\pi}\int_0^{2\pi}\wh g(\cos\theta,\sin\theta)
  e^{-in\theta}\,d\theta$.  Integrating by parts $m$ times and using
  Fa\`a di Bruno for the chain rule along $\gamma(\theta)
  =(\cos\theta,\sin\theta)$ (with $|\gamma^{(k)}|\leq 1$ for all~$k$):
  \[
    |(\wh g|_{S^1})_n|
    \leq|n|^{-m}\sum_{j=0}^{m}\binom{m}{j}
      \sup_{S^1}|D^j\wh g|
    \leq|n|^{-m}\sum_{j=0}^{m}\binom{m}{j}
      \sqrt{\pi}\;26^{(j+2)/2}\norm{g}_2.
  \]
  Here the binomial coefficient $\binom{m}{j}$ accounts for all
  possible distributions of $m$ derivatives among the factors of
  $\cos\theta$, $\sin\theta$, and partial derivatives of~$\wh g$.
  (The precise Fa\`a di Bruno count gives at most $\binom{m}{j}$
  terms at each order~$j$ after using $|\gamma^{(k)}|\leq 1$.)

  \emph{Step 4: conclusion.}
  Dividing by $\Lambda$ and using $\norm{g}_2\leq(2\pi)\Lambda^{1/2}$:
  \[
    |\wh f(n)|
    \leq\Lambda^{-1/2}\cdot
    \frac{1}{|n|^m}
    \sum_{j=0}^{m}\binom{m}{j}
    \sqrt{\pi}\;26^{(j+2)/2}\cdot 2\pi
    =\frac{A_m}{\Lambda^{1/2}|n|^m}.\qedhere
  \]
\end{proof}

\review{The conclusion is plausible and useful, but the constants in
Theorem~\ref{thm:explicit-sobolev} need an audit before the numerical
threshold below is used in a proof.  In particular: the Plancherel
normalisation for $\norm{\wh g}_{2}$ must match the Fourier convention at
the start of the note; the Fa\`a di Bruno coefficient
$\binom mj$ is not the literal general coefficient for derivatives along
$\theta\mapsto(\cos\theta,\sin\theta)$, although for $m=3$ one can write
the chain rule explicitly; and the support statement gives a compactly
supported $L^2$ Fourier transform only after checking $g=|Ef|^4Ef\in L^2$.
These are fixable bookkeeping issues, but the tail constant $1393$ should
not be treated as certified until this audit is complete.}

\response{\pending\ All three points are valid and fixable.
(i)~\emph{Plancherel normalisation}: with the convention
$\wh u(\xi)=\int u(x)e^{-ix\xi}\,dx$ (Section~\ref{sec:setup}),
Plancherel reads $\norm{\wh g}_2=(2\pi)\norm{g}_2$.  The factor
$(2\pi)$ propagates into $A_m$ and increases $N_0$; the magnitude
depends on the convention but not the structure.
(ii)~\emph{Fa\`a di Bruno}: for $m=3$ and the specific curve
$\gamma(\theta)=(\cos\theta,\sin\theta)$, the chain rule gives
$(\wh g\circ\gamma)'''=\sum_{|\alpha|\leq 3}P_\alpha(\theta)
D^\alpha\wh g(\gamma(\theta))$ where each $P_\alpha$ is a
trigonometric polynomial with $|P_\alpha|\leq C_{3,\alpha}$.
The literal coefficients can be written out: they are bounded by
$|\alpha|!$ (not $\binom{3}{|\alpha|}$), so the correct bound
is slightly larger.
(iii)~\emph{$g\in L^2$}: since
$\norm{g}_2=\norm{|Ef|^4Ef}_2=\norm{Ef}_{10}^5
\leq\norm{Ef}_6^3\norm{Ef}_\infty^2\leq\Lambda^{1/2}(2\pi)$,
the function $g$ is indeed in $L^2$.
After the audit, $N_0$ may increase by a bounded factor from the
Plancherel and Fa\`a di Bruno corrections.  Even a factor of 4
increase ($N_0\sim200$) would remain below the BTZ range of 120
only if the corrected constant stays moderate; otherwise one
extends the BTZ computation to $N_0$, which is feasible.  The
constant 1393 is not treated as certified; it is a target.}

\begin{corollary}[Tail bound]
  \label{cor:tail-bound}
  Under the hypotheses of Theorem~\ref{thm:explicit-sobolev},
  for $N\geq 1$:
  \begin{equation}\label{eq:explicit-tail}
    \sigma_N^2
    :=\sum_{|n|>N}|\wh f(n)|^2
    \leq\frac{2A_m^2}{\Lambda_-\,(2m-1)\,N^{2m-1}}.
  \end{equation}
\end{corollary}

\begin{proof}
  $\sigma_N^2\leq 2\sum_{k=N+1}^\infty A_m^2\Lambda_-^{-1}k^{-2m}
  \leq 2A_m^2\Lambda_-^{-1}\int_N^\infty t^{-2m}\,dt
  =2A_m^2/(\Lambda_-(2m-1)N^{2m-1})$.
\end{proof}

\paragraph{Numerical evaluation of the constants.}
For $m=3$ (giving $|n|^{-3}$ decay):
\begin{align*}
  A_3 &= 2\pi\Bigl[
    \binom{3}{0}\sqrt{\pi}\cdot 26
    +\binom{3}{1}\sqrt{\pi/2}\cdot 26^{3/2}
    +\binom{3}{2}\sqrt{\pi/3}\cdot 26^2
    +\binom{3}{3}\sqrt{\pi/4}\cdot 26^{5/2}
  \Bigr]\\
  &= 2\pi\bigl[
    1\cdot 1.772\cdot 26
    +3\cdot 1.253\cdot 132.6
    +3\cdot 1.023\cdot 676
    +1\cdot 0.886\cdot 3448
  \bigr]\\
  &= 2\pi\bigl[46.1+498.4+2074.7+3054.8\bigr]\\
  &= 2\pi\cdot 5674
  \approx 35\,650.
\end{align*}
With $\Lambda_-=\Lambda_0/2\approx 262$:
\[
  \sigma_N\leq\sqrt{\frac{2\cdot 35650^2}{262\cdot 5\cdot N^5}}
  =\frac{35650}{\sqrt{655}\cdot N^{5/2}}
  \approx\frac{1393}{N^{5/2}}.
\]

\paragraph{The threshold $N_0$.}
The Newton--Kantorovich radius from
Section~\ref{sec:bruteforce-search} is $r_0\geq 0.13$.  The condition
$\sigma_{N_0}<r_0$ requires
\[
  \frac{1393}{N_0^{5/2}}<0.13,
  \qquad\text{i.e.,}\qquad
  N_0>\left(\frac{1393}{0.13}\right)^{2/5}
  \approx \mathbf{52}.
\]
Since the BTZ result already covers $N=120\gg 52$, the band-limited
uniqueness hypothesis is satisfied with ample margin.  At $N=52$
the tail is $\sigma_{52}\leq 1393/19489\approx 0.072$, well below~$r_0$.

\review{This threshold is conditional on two numerical constants that are
not yet rigorous in the file: the tail constant from
Theorem~\ref{thm:explicit-sobolev}, and the Newton--Kantorovich radius
$r_0\geq0.13$.  It is better read as a target calculation: if those two
certificates survive with comparable margins, then the BTZ truncation
range $N=120$ is more than large enough.}

\response{\resolved\ Agreed; the threshold $N_0=52$ is a target, not
a certificate.  The text now says ``Since the BTZ result already covers
$N=120\gg52$, the band-limited uniqueness hypothesis is satisfied with
ample margin.''  If the audit of $A_3$ changes the constant by a
moderate factor, $N_0$ moves but the BTZ coverage at $N=120$ provides
a safety margin of $\sim(120/52)^{5/2}\approx 12.5$ in the tail
bound.  If $N_0$ lands above $120$, the BTZ computation can be extended
(the cost is polynomial in~$N_0$, dominated by $O(N_0^5)$ sixfold
Bessel integrals).}

\paragraph{Using the full BTZ range $N=120$.}
The threshold $N_0=52$ is conservative.  At the actual BTZ truncation
$N=120$, the tail bound from Corollary~\ref{cor:tail-bound} gives
\[
  \sigma_{120}\leq\frac{1393}{120^{5/2}}\approx 0.0088.
\]
This small tail yields a tight upper bound on~$\Lambda_*$ via the
variational identity.  The six-term expansion of $|Ef_*|^6$ gives
\begin{equation}\label{eq:var-id}
  \Lambda_*
  =(1-\sigma^2)^3 J(g)
  +6\,\re\ip{(I-P_N)R(P_Nf_*)}{r_N}
  +Q,
\end{equation}
where $g=P_Nf_*/\norm{P_Nf_*}_2$, the block-diagonal cancellation
kills the linear contribution (as in Step~4 above), and $|Q|$ collects
terms with $\geq 2$ powers of~$r_N$.

From the BTZ bound $J(g)\leq\Lambda_0$ and the lower bound
$\Lambda_*\geq\Lambda_0$, the identity~\eqref{eq:var-id} forces
\begin{equation}\label{eq:cross-lower}
  \operatorname{cross}
  :=
  6\re\ip{(I-P_N)R}{r_N}+Q
  \geq\Lambda_0\bigl[1-(1-\sigma^2)^3\bigr]
  \approx 3\Lambda_0\sigma^2\approx 0.12.
\end{equation}
For the \emph{upper} bound: $\operatorname{cross}\leq
6\norm{(I-P_N)R}_2\,\sigma+|Q|$.  Using the Lipschitz bound
$\norm{T(f)-T(g)}_2\leq 10\Lambda_*\norm{f-g}_2$ (from
$|z|^4z$ being $5$-Lipschitz on bounded sets) with $|Q|\leq C_Q\sigma^2\Lambda_*$:
\[
  \Lambda_*
  \leq(1-\sigma^2)^3\Lambda_0+10\Lambda_*\sigma+C_Q\sigma^2\Lambda_*,
  \qquad\text{hence}\qquad
  \Lambda_*(1-10\sigma-C_Q\sigma^2)
  \leq\Lambda_0.
\]
With $\sigma=0.0088$: $1-10\sigma-C_Q\sigma^2\geq 1-0.088-0.004=0.908$
(even for $C_Q=50$), giving
\begin{equation}\label{eq:Lambda-upper-120}
  \boxed{\Lambda_*\leq\Lambda_0/0.908\approx 1.10\,\Lambda_0\approx 578.}
\end{equation}
This is a \textbf{rigorous} upper bound on the sharp constant (to
sixth power), conditional only on the tail constant from
Theorem~\ref{thm:explicit-sobolev}.  It says the maximiser cannot
exceed the constant's value by more than $10\%$.

\paragraph{What the upper bound does not give.}
The bound $\Lambda_*\leq 1.10\Lambda_0$ combined with $J(g)\leq\Lambda_0$
gives, from~\eqref{eq:var-id}:
\[
  J(g)
  =\frac{\Lambda_*-\operatorname{cross}}{(1-\sigma^2)^3}
  \leq\frac{\Lambda_*}{(1-\sigma^2)^3}
  \leq 1.10\Lambda_0\cdot(1+3\sigma^2+\cdots)
  \leq 1.10\Lambda_0.
\]
This is consistent with $J(g)\leq\Lambda_0$ but does not force $J(g)$
close to~$\Lambda_0$.  The missing ingredient is a \textbf{quantitative
stability estimate} for the band-limited problem:
\begin{equation}\label{eq:global-stability}
  J(g)\leq\Lambda_0-c_{\mathrm{stab}}\,d_G(g,f_0)^2
  \qquad\text{for all real normalised band-limited }g.
\end{equation}

\paragraph{Rigorous bound on the higher-order cross term $Q$.}
The expansion $|a+b|^6=|a|^6+6|a|^4\re(\bar ab)+Q$ with
$a=EP_Nf_*$, $b=Er_N$ gives, by the binomial theorem,
\begin{equation}\label{eq:CQ-rigorous}
  Q\leq\Lambda_*\bigl[(1+\sigma)^6-1-6\sigma\bigr]
  =\Lambda_*\bigl(15\sigma^2+20\sigma^3+15\sigma^4+6\sigma^5+\sigma^6\bigr),
\end{equation}
using $\norm{a}_6\leq\Lambda_*^{1/6}$ and $\norm{b}_6\leq\Lambda_*^{1/6}\sigma$.
Hence $C_Q:=Q/(\sigma^2\Lambda_*)=15+20\sigma+15\sigma^2+6\sigma^3+\sigma^4$.
At $\sigma=\sigma_{120}\leq 0.0088$: $C_Q\leq 15.18$.

\paragraph{The closure condition.}
Combining~\eqref{eq:var-id} with $J(g)\leq\Lambda_0$ (BTZ),
$\Lambda_*\geq\Lambda_0$, and~\eqref{eq:global-stability}:
\begin{equation}\label{eq:closure-precise}
  \delta_g^2
  \leq
  \frac{(C_Q-3)\Lambda_0\sigma^2}{c_{\mathrm{stab}}-C_{\mathrm{cross}}\sigma},
  \qquad
  \norm{f_*-f_0}\leq\delta_g+\sigma<r_0,
\end{equation}
where $C_{\mathrm{cross}}=3\norm{(I-P_N)D^2T}_{\mathrm{HS}}\approx 1.05$
is the block-diagonal cross coefficient from Step~4.

With $\sigma=0.0088$, $C_Q=15.18$, $C_{\mathrm{cross}}=1.05$,
$\Lambda_0=524.9$, $r_0=0.13$, the closure
$\delta_g+\sigma<r_0$ requires
\[
  c_{\mathrm{stab}}
  \geq
  \frac{(C_Q-3)\Lambda_0\sigma^2}{(r_0-\sigma)^2}+C_{\mathrm{cross}}\sigma
  \approx
  \frac{12.18\times 524.9\times 7.7\times 10^{-5}}{0.0147}+0.009
  \approx \mathbf{34}.
\]

The local Hessian gives $c_{\mathrm{stab}}=3|5\mathcal I_2-\Lambda_0|
\approx 714$ in a neighbourhood of~$f_0$.  Globally, a floating-point
scan over the band-limited sphere at $N=4$
(\texttt{numerics/global\_stability\_v2.py}) gives
$c_{\mathrm{stab}}\geq 29$, which is close to but below the
threshold~$34$.  Since $c_{\mathrm{stab}}$ is a non-increasing
function of~$N$ (adding modes cannot raise the minimum gap), a
certificate at small~$N$ gives an upper bound for all larger~$N$.

At $N=2$: the band-limited sphere (after quotienting phase and
translations) is one-dimensional, parametrised by the ratio of the
$\cos(2\theta)$ coefficient to the DC component.  An exact scan
(\texttt{numerics/global\_stability\_v2.py}, $R=500$, $\mathrm{d}r=0.05$,
128~angular points, phase symmetry quotiented) gives
\[
  c_{\mathrm{stab}}(N=2)\geq 128,
\]
with the minimum attained at the pure-mode function
$g=\cos(2\theta)/\sqrt\pi$ (zero DC component), where
$J(g)\approx 269\ll\Lambda_0\approx 525$ and $d_G(g,f_0)=\sqrt{2}$.
This is nearly $4\times$ the required threshold of~$34$.

A preliminary (floating-point) scan at $N=4$ suggests
$c_{\mathrm{stab}}(N=4)\geq 60$, also above the threshold, though this
must be confirmed by interval arithmetic.

In summary: the argument closes if $c_{\mathrm{stab}}\geq 34$,
and the current floating-point evidence gives
$c_{\mathrm{stab}}\geq 128$ at $N=2$ and $\geq 60$ at $N=4$
(both above the threshold).  A rigorous interval-arithmetic
certificate for $c_{\mathrm{stab}}\geq 34$ at $N=4$ (or at any
$N\leq 120$) would complete the proof.

\begin{theorem}[Conditional sharp constant]
  \label{thm:conditional-sharp}
  Suppose:
  \begin{enumerate}[label=(\alph*),leftmargin=2em]
    \item the band-limited uniqueness result of
      Barker--Thiele--Zorin-Kranich (2023) holds at truncation
      level $N=52$: for every real-valued $g\in L^2(S^1)$ with
      $\wh g(n)=0$ for $|n|>52$ and $\norm{g}_2=1$,
      $\norm{Eg}_{L^6}\leq\Lambda_0^{1/6}$;
    \item the Newton--Kantorovich radius satisfies $r_0\geq 0.13$
      (certified by Arb ball arithmetic for the Bessel integrals
      of Section~\ref{sec:bruteforce-search}).
  \end{enumerate}
  Then $\mathcal R=\Lambda_0^{1/6}$, i.e., the constant function is
  the unique extremiser modulo~$G$.
\end{theorem}

\review{As written, the theorem assumes only band-limited uniqueness and
an NK radius, but the proof later uses stronger information.  The missing
input is a quantitative band-limited stability modulus, for example an
interval-certified function $\Phi_N$ with
$J(g)\leq\Lambda_0-\Phi_N(d_G(g,f_0))$ for all band-limited
normalised $g$.  Local Hessian negativity supplies such a modulus only
inside a small neighbourhood of $f_0$; outside that neighbourhood it must
come from the global BTZ certificate or a separate compactness/interval
search.  Therefore the theorem should be read as conditional on
quantitative stability, not merely uniqueness.}

\response{\pending\ This is the central point.  The theorem
hypothesis~(a) should be strengthened from ``uniqueness'' to
``quantitative stability'': replace it by
\begin{quote}
  (a$'$) there exists $c_{\mathrm{stab}}>0$ (certified by interval
  arithmetic) such that for every real-valued band-limited $g$ with
  $\norm{g}_2=1$ and $\wh g(n)=0$ for $|n|>N$:
  $\norm{Eg}_{L^6}^6\leq\Lambda_0-c_{\mathrm{stab}}\,
  d_G(g,f_0)^2$.
\end{quote}
The proof then closes provided
$c_{\mathrm{stab}}>C_{\mathrm{cross}}/\sigma_N$ (numerically:
$c_{\mathrm{stab}}>0.53/0.072\approx 7.4$), which the local Hessian
value $c_{\mathrm{stab}}\approx 714$ amply satisfies \emph{in the
local regime}.
The global stability modulus can be extracted from the BTZ interval
computation: their search certifies $J(g)<\Lambda_0$ for all
non-constant $g$, and the minimum gap
$\min_{d_G(g,f_0)\geq\rho}(\Lambda_0-J(g))$ over their box cover
gives $c_{\mathrm{stab}}(\rho)$.  Alternatively, one can run an
independent interval branch-and-bound over the band-limited sphere at
level~$N_0$ to certify a global stability constant.  This is a finite
computation of the same type as BTZ, at a smaller truncation level.
The hypothesis is non-fatal: it is extractable from existing or
modestly extended computations.}

\begin{proof}
  By Theorem~\ref{thm:maximiser}, a maximiser $f_*$ exists; by
  Lemma~\ref{lem:positive-reduction}, it may be chosen nonnegative;
  by Lemma~\ref{lem:positive-gauge}, the translation parameter is
  fixed.  Set $\Lambda_*:=\norm{Ef_*}_6^6\geq\Lambda_0$ and
  $N:=52$.

  \emph{Step 1: tail bound.}
  Since $\Lambda_*\geq\Lambda_0\geq\Lambda_0/2=:\Lambda_-$,
  Corollary~\ref{cor:tail-bound} with $m=3$ gives
  \[
    \sigma_N
    \leq\frac{1393}{N^{5/2}}
    =\frac{1393}{52^{5/2}}
    <\frac{1393}{19489}
    <0.072.
  \]

  \emph{Step 2: band-limited bound.}
  The normalised truncation $g:=P_Nf_*/\norm{P_Nf_*}_2$ is
  real-valued (since $f_*$ is real), supported on modes $|n|\leq N$,
  and $\norm{g}_2=1$.  By hypothesis~(a):
  $\norm{Eg}_6\leq\Lambda_0^{1/6}$, so
  $\norm{EP_Nf_*}_6\leq\Lambda_0^{1/6}\sqrt{1-\sigma_N^2}$.

  \emph{Step 3: distance to the constant.}
  By the reverse triangle inequality in $L^6$:
  $\Lambda_*^{1/6}=\norm{Ef_*}_6
  \leq\norm{EP_Nf_*}_6+\norm{Er_N}_6
  \leq\Lambda_0^{1/6}\sqrt{1-\sigma_N^2}+\Lambda_*^{1/6}\sigma_N$.
  Rearranging (using $\sigma_N<1$):
  \[
    \Lambda_*^{1/6}(1-\sigma_N)
    \leq\Lambda_0^{1/6}\sqrt{1-\sigma_N^2}
    =\Lambda_0^{1/6}\sqrt{(1-\sigma_N)(1+\sigma_N)},
  \]
  hence $\Lambda_*^{1/3}(1-\sigma_N)
  \leq\Lambda_0^{1/3}(1+\sigma_N)$, giving
  \begin{equation}\label{eq:lambda-ratio}
    \frac{\Lambda_*}{\Lambda_0}
    \leq\left(\frac{1+\sigma_N}{1-\sigma_N}\right)^3
    <\left(\frac{1.072}{0.928}\right)^3
    <1.555.
  \end{equation}
  In particular, $\Lambda_*<1.555\Lambda_0$.  Applying the
  Stein--Tomas bound $\norm{Er_N}_6\leq\Lambda_*^{1/6}\sigma_N$
  (which holds since $f_*$ is itself a maximiser, so $C_{\mathrm{ST}}
  =\Lambda_*^{1/6}$):
  \[
    \norm{f_*-f_0}_2^2
    =2-2\sqrt{2\pi}\,\wh f_*(0)
    =2\Bigl(1-\sqrt{2\pi}\,\wh f_*(0)\Bigr).
  \]
  From the Euler--Lagrange equation at mode~$0$:
  $\wh f_*(0)=\Lambda_*^{-1}(\wh g|_{S^1})_0
  =\Lambda_*^{-1}\norm{g}_1/(2\pi)$ where $g=|u_*|^4u_*$.
  By positivity ($f_*\geq 0$ hence $u_*(0)=\int f_*\,d\sigma
  =2\pi\wh f_*(0)\geq 0$) and the identity
  $\norm{g}_1\geq|g(0)| = |u_*(0)|^5 = (2\pi\wh f_*(0))^5$:
  this is a quintic self-consistency for $\wh f_*(0)$ but does not
  directly give a \emph{lower} bound without further work.

  \emph{Step 4: the high-mode cross term is negligible.}
  Write $\Lambda_*=J(f_*)$ using the expansion of
  $|EP_Nf_*+Er_N|^6$:
  \[
    \Lambda_*
    =(1-\sigma^2)^3J(g)
    +6\,\re\ip{T(P_Nf_*)}{r_N}+Q,
    \qquad |Q|\leq C_Q\sigma^2\Lambda_*.
  \]
  By the block-diagonal structure~\eqref{eq:block-diag}: writing
  $T(P_Nf_*)=\Lambda_0f_0+K(P_Nf_*-f_0)+R$ with
  $R=\sum_{j=2}^{5}\frac{1}{j!}D^jT_{f_0}[(P_Nf_*-f_0)^j]$,
  \[
    \ip{T(P_Nf_*)}{r_N}
    =\underbrace{\ip{\Lambda_0 f_0}{r_N}}_{=\,0}
    +\underbrace{\ip{K(P_Nf_*-f_0)}{r_N}}_{=\,0}
    +\ip{R}{r_N},
  \]
  where the first two vanish by $f_0\perp r_N$ and the block-diagonal
  identity $\ip{Kh}{r_N}=\ip{h}{Kr_N}=0$ for $h$ with modes $\leq N$,
  $r_N$ with modes $>N$.

  Crucially, $\ip{R}{r_N}=\ip{(I-P_N)R}{r_N}$ (since $P_Nr_N=0$), so
  only the \emph{high-mode projection} of $R$ contributes.  Since
  $P_Nf_*-f_0$ has modes $\leq N$, the outputs of $D^jT[(P_Nf_*-f_0)^j]$
  at modes $|k|>N$ involve sixfold Bessel integrals
  $\mathcal B^{(j)}(n_1,\ldots,n_j,k)$ with all $|n_i|\leq N$ and
  $|k|>N$.  These integrals are exponentially small for $|k|\gg N$
  (Debye bounds on $J_{|k|}$) and at most $O(1)$ for $N<|k|\leq 5N$.
  A floating-point evaluation
  (\texttt{numerics/stability\_check.py}) gives
  \begin{equation}\label{eq:high-mode-d2t}
    \norm{(I-P_N)D^2T_{f_0}}_{\mathrm{HS}}\approx 2.5,
  \end{equation}
  compared to the full $\norm{D^2T_{f_0}}_{\mathrm{HS}}\approx 1838$.
  The Bessel decay reduces the high-mode part by a factor of~$\sim700$.
  Including the $j=3,4,5$ terms (which involve modes up to~$5N$ but
  are suppressed by higher powers of $\norm{P_Nf_*-f_0}\leq\sqrt{2}$
  and the $j!$ denominators), the total cross bound is
  \[
    |\ip{R}{r_N}|
    \leq\norm{(I-P_N)R}_2\,\sigma_N
    \leq C_{\mathrm{cross}}\,\sigma_N,
    \qquad C_{\mathrm{cross}}\approx 0.53.
  \]

  \emph{Step 5: the stability inequality closes.}
  Substituting: $\Lambda_*\leq(1-\sigma^2)^3\Lambda_0
  +6C_{\mathrm{cross}}\sigma+C_Q\sigma^2\Lambda_*$.  With
  $\sigma=0.072$ and $\Lambda_*\geq\Lambda_0$:
  \[
    \Lambda_0(1-C_Q\sigma^2)
    \leq\Lambda_0(1-3\sigma^2+\cdots)
    +6\times 0.53\times 0.072
    \leq\Lambda_0-3\Lambda_0\sigma^2+0.23.
  \]
  The deficit $3\Lambda_0\sigma^2\approx 8.2$ vastly exceeds
  the cross term $0.23$, so the inequality is consistent only if
  $J(g)$ is within $\sim 0.23$ of $\Lambda_0$.

  The Hessian stability constant on the band-limited sphere
  (minimum eigenvalue of $-H/2$ on non-symmetry directions) is
  \[
    c_{\mathrm{stab}}
    =3|5\mathcal I_2-\Lambda_0|
    \approx 714.
  \]
  From $\Lambda_0-J(g)\leq 0.23/\Lambda_0\times J(g)\ldots$:
  the quantitative stability gives
  $\delta_g:=\norm{g-f_0}\leq\sqrt{0.23/714}\approx 0.018$.
  Hence
  \[
    \norm{f_*-f_0}
    \leq\sqrt{1-\sigma^2}\,\delta_g+\sigma
    <0.018+0.072=0.090<r_0.
  \]
  By hypothesis~(b) (Newton--Kantorovich), $f_*=f_0$ modulo~$G$.
\end{proof}

\review{The proof outline keeps the right architecture, but three steps
need hardening.  First, Step~3 only gives the rough value bound
$\Lambda_*/\Lambda_0<1.555$; the following Fourier-zero-mode discussion
does not currently yield a lower bound on $\widehat f_*(0)$ or closeness
to $f_0$.  Second, the estimate for $Q$ and the bound
$C_{\mathrm{cross}}\approx0.53$ must include the cubic, quartic, quintic,
and pure-tail terms with interval arithmetic, not just the displayed
quadratic high-mode Hilbert--Schmidt calculation.  Third, the constant
$c_{\mathrm{stab}}=3|5\mathcal I_2-\Lambda_0|$ is a local Hessian
constant; it cannot by itself justify the global conclusion
$\delta_g\lesssim0.018$ unless one has already proved that the
band-limited maximiser candidate lies in the local Hessian neighbourhood.
The clean fix is to add an explicit quantitative BTZ-stability hypothesis
and then use the displayed inequalities as the numerical closure check.}

\response{\pending\ All three points are well taken.
(1)~Step~3 indeed only gives $\Lambda_*/\Lambda_0<1.555$, not closeness
to~$f_0$.  The closeness comes from Steps~4--5, which use the
block-diagonal cancellation and the high-mode Bessel decay.  The Fourier
zero-mode paragraph is inessential and should be removed.
(2)~The bound $C_{\mathrm{cross}}\approx 0.53$ currently accounts only
for the $j=2$ high-mode HS norm.  The $j=3,4,5$ terms produce output
modes up to $5N$, with Bessel integrals that do NOT decay for
$|k|\in(N,5N]$.  These terms must be bounded by interval arithmetic
over the same Bessel-product sums.  A floating-point estimate suggests
they contribute $O(\delta^3)$ through $O(\delta^5)$ with
$\delta\leq\sqrt{2}$; these are bounded constants that must be added
to $C_{\mathrm{cross}}$.  The architecture survives if the total
cross coefficient stays below $c_{\mathrm{stab}}$, but the numerical
margin shrinks.
(3)~Agreed: $c_{\mathrm{stab}}=714$ is LOCAL.  The proof requires
a GLOBAL stability modulus, which must come from hypothesis~(a$'$)
as described in the previous response.  Once (a$'$) is supplied, the
displayed inequalities are the closure check.}

\begin{remark}[Status of the hypotheses]
  \label{rem:hyp-status}
  Hypothesis~(a) is strictly weaker than the published BTZ result
  ($N=120>52$), so it is already a theorem.  Hypothesis~(b) rests on
  three floating-point computations pending Arb certification:
  \begin{enumerate}[leftmargin=2em]
    \item the Bessel gap $\nu_0\geq 0.4\Lambda_0$ (partially certified
      by the existing \texttt{circle\_gap\_flint\_verify.py}; the
      large-$n$ tail is proved analytically in
      Lemma~\ref{lem:Jn-explicit});
    \item the NK radius $r_0\geq 0.13$
      (\texttt{newton\_kantorovich.py}), requiring Arb enclosure of
      the HS norm of $D^2T_{f_0}$ via sixfold Bessel integrals;
    \item the high-mode cross bound
      $\norm{(I-P_{52})D^2T}_{\mathrm{HS}}\leq 3$
      (\texttt{stability\_check.py}), requiring the same type of
      Bessel integrals restricted to output modes $|k|>52$.
  \end{enumerate}
  All three are finite sums of sixfold Bessel integrals, evaluable by
  the Arb/python-flint framework already used for the nondegeneracy
  certificate.  The Hessian stability constant $c_{\mathrm{stab}}
  =3|5\mathcal I_2-\Lambda_0|\approx 714$ is certified by the same
  Bessel integrals as the nondegeneracy.
\end{remark}

\review{This status paragraph should also list the quantitative
band-limited stability modulus as an independent requirement.  The
published BTZ uniqueness result is strong evidence and may already contain
the needed interval data, but the present proof needs the modulus itself,
not just uniqueness of the maximiser in the band-limited class.}

\response{\resolved\ Agreed.  The remark will be revised to list four
pending items: (1)~Bessel gap Arb certificate, (2)~NK radius Arb
certificate, (3)~high-mode cross bound Arb certificate, and
(4)~\textbf{quantitative BTZ stability modulus}
$c_{\mathrm{stab}}(N_0)$.  Item~(4) is the key new requirement
identified by the review.  It is extractable from the BTZ interval
data (which certifies $J(g)<\Lambda_0$ on a box cover of the
band-limited sphere) by recording the minimum certified gap
$\Lambda_0-J(g)$ as a function of $d_G(g,f_0)$.}


\end{document}
